\section{Questão 11}
Considere que uma variável randômica $X$ que diz respeito aos resultados de um
determinado exame médico para um paciente que está ou
não com determinada enfermidade.
A tabela abaixo registra algumas amostras dos pares $(X=x_i, C_i)$, onde $C_i=1$ se o paciente possui a enfermidade e $C_i=0$ caso contrário. 

\begin{table}[htbp]
    \centering
    \begin{tabular}{c | ccc ccc ccc ccc}
    \toprule
        $x_i$   & 1.5 & 2.8 & 1.2 & 2.2 & 3.8 & 5.4 & 5.1 & 2.8 & 4.4 & 2.3 & 4.5 & 3.8 \\
    \midrule
        $C_i$ & 0   & 0   & 1   & 0   & 1   & 1   & 1   & 0   & 1   & 0   & 1   & 1 \\
    \bottomrule
    \end{tabular}
\end{table}

Considere que o conjunto de dados mensurados possam ser descritos por uma PDF Gaussiana. 
Usando a Lei de Bayes determine a probabilidade dos resultados dos exames serem classificados como enfermo ou
sadio: $X = 5.7$; $X = 3.0$; $X = 2.6$.

\respostas

O classificador atribui a uma amostra $x_i$ a classe $C_i$ que otimize a probabilidade condicional 
$P(C_i | x_i)$. Pela Lei de Bayes, essa condicional (a distribuição \textit{a posteriori}) 
é proporcional ao produto da \textit{priori} pela \textit{verossimilhança}:
%
\begin{equation}
    P(C_i = k | x_i) 
    = \frac{P(C_i = k) P(x_i | C_i)}{P(x_i)}
    = \frac{P(C_i) P(x_i | C_i)}{\sum_k P(C_i=k)P(x_i | C_i=k)}
    \propto P(C_i) P(x_i | C_i)
\end{equation}
%
e portanto o classificador Bayesiano busca encontrar a classe que maximize $P(C_i) P(x_i | C_i)$.
A probabilidade \textit{a priori} de cada classe (com ou sem enfermidade) pode ser estimada diretamente da tabela:
%
\begin{align}
    P(C_i=0) &\approx \frac{N_0}{N_1 + N_0} = \frac{5}{12}, &
    P(C_i=1) &\approx \frac{N_1}{N_1 + N_0} = \frac{7}{12},
\end{align}
%
com 
$N_0=5$ a quantidade de pacientes sem a enfermidade, e 
$N_1=7$ a quantidade de pacientes com a enfermidade.
O cálculo das verossimilhanças requer mais suposições;
como é garantido que as classes vêm de PDFs Gaussianas, tem-se
%
\begin{equation}
    P(x | C_i = k) = 
    \frac{1}{\sqrt{2\pi\sigma_k^2}}
    \exp\left(
        -\frac{(x-\mu_k)^2}{2\sigma_k^2}
    \right),
\end{equation}
%
com $\mu_k$, $\sigma_k^2$ a média e variância da classe $k$.
Voltando ao conjunto de dados na tabela, tem-se
%
\begin{align}
    \mu_0 &\approx m_0 = \frac{1}{N_0} \sum_{x_i\in C_0} x_i = \frac{11.6}{5} = 2.320, &
    \sigma_0^2 &\approx \frac{1}{N_0-1} \sum_{x_i\in C_0} (x_i - m_0)^2 = 0.287, \\
    \mu_1 &\approx m_1 = \frac{1}{N_1} \sum_{x_i\in C_1} x_i = \frac{28.2}{7} = 4.029, &
    \sigma_1^2 &\approx \frac{1}{N_1-1} \sum_{x_i\in C_1} (x_i - m_1)^2 = 1.916.
\end{align}
%
Substituindo os momentos e aplicando a lei de Bayes para cada valor de teste, 
as probabilidades de enfermidade são encontradas.

\begin{table}[htbp!]
    \centering
    \begin{tabular}{c c c}
    \toprule
        $x_i$ & $P(C_0 | x_i)$ [\%] & $P(C_1 | x_i)$ [\%] \\
    \midrule
        5.7 & $<10^{-6}$ & $99.99$ \\
        3.0 & $52.08$ & $47.92$ \\
        2.6 & $73.28$ & $26.72$ \\
    \bottomrule
    \end{tabular}
\end{table}